% ---- pg-325 （书页 313） ----

均方误差在上式右端最后两项为零时达到极小, 这可通过先取 $\mathbf{b}=\boldsymbol{\Sigma}_{ZZ}^{-1}\boldsymbol{\sigma}_{ZY}=\boldsymbol{\beta}$, 然后取 $b_0=\mu_Y-(\boldsymbol{\Sigma}_{ZZ}^{-1}\boldsymbol{\sigma}_{ZY})'\boldsymbol{\mu}_Z=\beta_0$ 来实现. 于是最小均方误差为 $\sigma_{YY}-\boldsymbol{\sigma}_{ZY}'\boldsymbol{\Sigma}_{ZZ}^{-1}\boldsymbol{\sigma}_{ZY}$.

其次, 注意到 $\operatorname{Cov}(b_0+\mathbf{b}'\mathbf{Z},Y)=\operatorname{Cov}(\mathbf{b}'\mathbf{Z},Y)=\mathbf{b}'\boldsymbol{\sigma}_{ZY}$, 故对所有 $b_0$ 和 $\mathbf{b}$ 有
\[
[\operatorname{Corr}(b_0+\mathbf{b}'\mathbf{Z},Y)]^2
=\frac{[\mathbf{b}'\boldsymbol{\sigma}_{ZY}]^2}{\sigma_{YY}(\mathbf{b}'\boldsymbol{\Sigma}_{ZZ}\mathbf{b})}
\]
应用推广的柯西-施瓦茨不等\hyperref[eq:2-49]{式(2-49)}, 并取 $\mathbf{B}=\boldsymbol{\Sigma}_{ZZ}$, 得到
\[
(\mathbf{b}'\boldsymbol{\sigma}_{ZY})^2\leqslant\mathbf{b}'\boldsymbol{\Sigma}_{ZZ}\mathbf{b}\,\boldsymbol{\sigma}_{ZY}'\boldsymbol{\Sigma}_{ZZ}^{-1}\boldsymbol{\sigma}_{ZY}
\]
或
\[
[\operatorname{Corr}(b_0+\mathbf{b}'\mathbf{Z},Y)]^2
\leqslant\frac{\boldsymbol{\sigma}_{ZY}'\boldsymbol{\Sigma}_{ZZ}^{-1}\boldsymbol{\sigma}_{ZY}}{\sigma_{YY}}
\]
其中等号对 $\mathbf{b}=\boldsymbol{\Sigma}_{ZZ}^{-1}\boldsymbol{\sigma}_{ZY}=\boldsymbol{\beta}$ 成立. 最大相关的另一种表达式由 $\boldsymbol{\sigma}_{ZY}'\boldsymbol{\Sigma}_{ZZ}^{-1}\boldsymbol{\sigma}_{ZY}=\boldsymbol{\sigma}_{ZY}'\boldsymbol{\beta}=\boldsymbol{\sigma}_{ZY}'\boldsymbol{\Sigma}_{ZZ}^{-1}\boldsymbol{\Sigma}_{ZZ}\boldsymbol{\beta}=\boldsymbol{\beta}'\boldsymbol{\Sigma}_{ZZ}\boldsymbol{\beta}$ 得到. \hfill$\blacksquare$

$Y$ 与其最佳线性预测量之间的相关系数称为总体多重相关系数
\phantomsection\label{eq:7-48}
\begin{equation}
\rho_{Y(\mathbf{Z})}=+\sqrt{\frac{\boldsymbol{\sigma}_{ZY}'\boldsymbol{\Sigma}_{ZZ}^{-1}\boldsymbol{\sigma}_{ZY}}{\sigma_{YY}}}
\tag{7-48}
\end{equation}
总体多重相关系数的平方 $\rho_{Y(\mathbf{Z})}^2$ 称为总体决定系数. 注意, 与两个变量的相关系数不同, 多重相关系数为正的平方根, 故 $0\leqslant\rho_{Y(\mathbf{Z})}\leqslant1$.

总体决定系数有一个重要解释. 根据结论 7.12, 用 $\beta_0+\boldsymbol{\beta}'\mathbf{Z}$ 预测 $Y$ 时的均方误差为
\phantomsection\label{eq:7-49}
\begin{equation}
\sigma_{YY}-\boldsymbol{\sigma}_{ZY}'\boldsymbol{\Sigma}_{ZZ}^{-1}\boldsymbol{\sigma}_{ZY}
=\sigma_{YY}-\sigma_{YY}\left(\frac{\boldsymbol{\sigma}_{ZY}'\boldsymbol{\Sigma}_{ZZ}^{-1}\boldsymbol{\sigma}_{ZY}}{\sigma_{YY}}\right)
=\sigma_{YY}(1-\rho_{Y(\mathbf{Z})}^2)
\tag{7-49}
\end{equation}
若 $\rho_{Y(\mathbf{Z})}^2=0$, 则 $\mathbf{Z}$ 没有预测 $Y$ 的能力. 另一个极端 $\rho_{Y(\mathbf{Z})}^2=1$ 意味着 $Y$ 可以无误差地被预测.

\medskip
\noindent\textbf{例 7.11}\quad (最佳线性预测量及其均方误差与多重相关系数的确定)

给定 $Y$, $Z_1$ 和 $Z_2$ 的均值向量和协方差矩阵
\[
\boldsymbol{\mu}=\left[\begin{array}{c}
\mu_Y\\
\hdashline
\boldsymbol{\mu}_Z
\end{array}\right]
=\begin{bmatrix}5\\2\\0\end{bmatrix}
\qquad\text{和}\qquad
\boldsymbol{\Sigma}=\left[\begin{array}{c:c}
\sigma_{YY}&\boldsymbol{\sigma}_{ZY}'\\
\hdashline
\boldsymbol{\sigma}_{ZY}&\boldsymbol{\Sigma}_{ZZ}
\end{array}\right]
=\left[\begin{array}{c:cc}
10&1&-1\\
\hdashline
1&7&3\\
-1&3&2
\end{array}\right]
\]
试确定(a)最佳线性预测量 $\beta_0+\beta_1Z_1+\beta_2Z_2$, (b)这个预测量的均方误差, (c)多重相关系数. 此外, 验证均方误差等于 $\sigma_{YY}(1-\rho_{Y(\mathbf{Z})}^2)$.

首先,
\[
\boldsymbol{\beta}=\boldsymbol{\Sigma}_{ZZ}^{-1}\boldsymbol{\sigma}_{ZY}
=\begin{bmatrix}7&3\\3&2\end{bmatrix}^{-1}\begin{bmatrix}1\\-1\end{bmatrix}
=\begin{bmatrix}0.4&-0.6\\-0.6&1.4\end{bmatrix}\begin{bmatrix}1\\-1\end{bmatrix}
=\begin{bmatrix}1\\-2\end{bmatrix}
\]
\[
\beta_0=\mu_Y-\boldsymbol{\beta}'\boldsymbol{\mu}_Z=5-[1,-2]\begin{bmatrix}2\\0\end{bmatrix}=3
\]
故最佳线性预测量为 $\beta_0+\boldsymbol{\beta}'\mathbf{Z}=3+Z_1-2Z_2$. 均方误差为
\[
\sigma_{YY}-\boldsymbol{\sigma}_{ZY}'\boldsymbol{\Sigma}_{ZZ}^{-1}\boldsymbol{\sigma}_{ZY}
=10-[1,-1]\begin{bmatrix}0.4&-0.6\\-0.6&1.4\end{bmatrix}\begin{bmatrix}1\\-1\end{bmatrix}
=10-3=7
\]
多重相关系数为
\[
\rho_{Y(\mathbf{Z})}=\sqrt{\frac{\boldsymbol{\sigma}_{ZY}'\boldsymbol{\Sigma}_{ZZ}^{-1}\boldsymbol{\sigma}_{ZY}}{\sigma_{YY}}}
=\sqrt{\frac{3}{10}}=0.548
\]
注意 $\sigma_{YY}(1-\rho_{Y(\mathbf{Z})}^2)=10\times(1-3/10)=7$ 等于均方误差. \hfill$\blacksquare$
